\chapter{移動と交通 — Transporte}
\unidadheader{19}{移動と交通}{Transporte}{Moverse por la ciudad y viajar}

\begin{tcolorbox}[vocabbox, title={📌 Vocabulario Nuevo}]
\begin{tabularx}{\linewidth}{llX}
\toprule
\textbf{Japonés} & \textbf{Español} & \textbf{Tipo} \\
\midrule
\vocabitem{電車}{でんしゃ}{tren eléctrico}{sust.}
\vocabitem{飛行機}{ひこうき}{avión}{sust.}
\vocabitem{自転車}{じてんしゃ}{bicicleta}{sust.}
\vocabitem{乗る}{のる}{subir / abordar}{v. gr.1}
\vocabitem{降りる}{おりる}{bajar / descender}{v. gr.2}
\vocabitem{乗り換える}{のりかえる}{hacer transbordo}{v. gr.2}
\vocabitem{切符}{きっぷ}{boleto / ticket}{sust.}
\vocabitem{〜番線}{ばんせん}{andén número ~}{sust.}
\vocabitem{終点}{しゅうてん}{última parada}{sust.}
\vocabitem{〜分かかる}{ふんかかる}{tardar ~ minutos}{expr.}
\bottomrule
\end{tabularx}
\end{tcolorbox}

\subsection*{🗣️ Frases Clave}

\frase{\ruby{新宿}{しんじゅく}まで\ruby{何}{なん}\ruby{分}{ぷん}かかりますか？}
  {¿Cuántos minutos se tarda hasta Shinjuku?}
\frase{\ruby{電車}{でんしゃ}で\ruby{約}{やく}20\ruby{分}{ぷん}かかります。}
  {En tren se tarda aproximadamente 20 minutos.}
\frase{〜\ruby{線}{せん}に\ruby{乗}{の}り\ruby{換}{か}えてください。}
  {Haga transbordo a la línea ~.}
\frase{〜\ruby{番線}{ばんせん}のホームです。}
  {Es el andén número ~.}

\subsection*{⚙️ Gramática}

\patron{Medio de transporte: で}
  {交通手段 + で + 行く/来る/帰る}
  {\ej{\ruby{自転車}{じてんしゃ}で\ruby{学校}{がっこう}に\ruby{行}{い}きます。}
    {Voy a la escuela en bicicleta.}
  \ej{バスで20\ruby{分}{ぷん}かかります。}
    {En autobús se tarda 20 minutos.}}

\patron{Subir / Bajar transporte: に乗る / を降りる}
  {〜に乗る (subir a ~) / 〜を降りる (bajar de ~)}
  {\ej{\ruby{山手線}{やまのてせん}に\ruby{乗}{の}ってください。}
    {Súbase a la línea Yamanote.}
  \ej{\ruby{渋谷}{しぶや}で\ruby{降}{お}りてください。}
    {Baje en Shibuya.}}

\begin{tcolorbox}[ejerciciobox, title={📝 Ejercicios Rápidos}]
\begin{enumerate}
  \item ¿Cómo preguntas cuánto tiempo tarda el autobús?
  \item Di cómo llegas a tu trabajo/escuela.
  \item Traduce: \ruby{次}{つぎ}の\ruby{駅}{えき}で\ruby{乗}{の}り\ruby{換}{か}えてください。
\end{enumerate}
\vspace{0.4em}
\textbf{Respuestas:}
1) バスで\ruby{何分}{なんぷん}かかりますか？\quad
2) (personal)\quad
3) Por favor, haga transbordo en la próxima estación.
\end{tcolorbox}

\begin{tcolorbox}[kanjibox, title={🀄 Kanjis de la Unidad}]
\begin{tabularx}{\linewidth}{cllXX}
\toprule
\textbf{Kanji} & \textbf{On} & \textbf{Kun} & \textbf{Significado} & \textbf{Vocab} \\
\midrule
\kanjirow{電}{でん}{—}{electricidad}{\ruby{電車}{でんしゃ} / \ruby{電話}{でんわ}}
\kanjirow{車}{しゃ}{くるま}{vehículo / coche}{\ruby{電車}{でんしゃ} / \ruby{車}{くるま}}
\kanjirow{飛}{ひ}{と}{volar}{\ruby{飛行機}{ひこうき} / \ruby{飛}{と}ぶ}
\kanjirow{機}{き}{はた}{máquina / avión}{\ruby{飛行機}{ひこうき} / \ruby{機会}{きかい}}
\kanjirow{乗}{じょう}{の}{subir / abordar}{\ruby{乗}{の}る / \ruby{乗車}{じょうしゃ}}
\bottomrule
\end{tabularx}
\end{tcolorbox}
